\section{Repository and How to Run the Code}

\begin{tabularx}{\textwidth}{@{}lL@{}}
  \textbf{GitHub username} & Graf-J \\
  \textbf{Repository} & \href{https://github.com/Graf-J/MLDS-Q4-490-Advanced-Algorithms-HW1}{\texttt{github.com/Graf-J/MLDS-Q4-490-Advanced-Algorithms-HW1}} (public)
\end{tabularx}

The project uses \href{https://docs.astral.sh/uv/}{uv} for Python~3.14 and its dependencies:

\begin{lstlisting}
git clone https://github.com/Graf-J/MLDS-Q4-490-Advanced-Algorithms-HW1.git
cd MLDS-Q4-490-Advanced-Algorithms-HW1
uv sync             # creates .venv with the exact versions from uv.lock
uv run jupyter lab  # or open the notebooks in VS Code with the .venv kernel
\end{lstlisting}

Every part of the assignment is one Jupyter notebook that runs from top to bottom. The dynamic-programming solvers of Section~\ref{sec:dp} live in \texttt{src/} and are shared by the notebooks; all figures are written to \texttt{report/figures/}.

\begin{center}
\small
\begin{tabularx}{\textwidth}{@{}l l L@{}}
  \toprule
  \textbf{Part} & \textbf{Notebook} & \textbf{Code and runtime} \\
  \midrule
  2: Environment & \texttt{2\_ValidateEnvironment.ipynb} & Environment in \texttt{src/environments/shuttleenvironment.py} \\
  2a: Enumeration & \texttt{2a\_Enumeration.ipynb} & Solver in \texttt{src/solvers/shuttleenumeration.py}; the solve takes about 6\,s \\
  2b: Value iteration & \texttt{2b\_ValueIteration.ipynb} & Solver in \texttt{src/solvers/shuttlevalueiteration.py} \\
  2c: Policy iteration & \texttt{2c\_PolicyIteration.ipynb} & Solver in \texttt{src/solvers/shuttlepolicyiteration.py} \\
  \midrule
  3a: CartPole & \texttt{3a\_CartPole-v1.ipynb} & Part~1: four trainings about 12\,min, seed experiment another 45\,min. Part~2 (baseline): training about 30\,min, checks about 3\,min; it only needs the definition cells of Part~1. \\
  3b: Pong & \texttt{3b\_Pong-v5.ipynb} & Part~1: training about 100\,min, evaluation and checks about 30\,min. Part~2 (baseline): training about 2\,h\,45\,min, evaluation and checks about 80\,min; it only needs the definition cells of Part~1. \\
  \bottomrule
\end{tabularx}
\end{center}

\begin{sloppypar}
\textbf{Pong without retraining.} The Pong runs are saved every 250 episodes to \texttt{checkpoints/3b/reinforce/} and \texttt{checkpoints/3b/actor\_critic/} (policy, critic and training log), so a crash loses at most 250 episodes. With the default flags \texttt{RETRAIN = False} and \texttt{RECOMPUTE = False}, the notebook loads the finished runs and the cached evaluations and checks instead of computing them again.
\end{sloppypar}

\begin{notebox}
\textbf{Validation tests.} Every validation test described in this report is implemented as a check in the \emph{Checks} section of the corresponding notebook. The checks run together with the notebook and print \texttt{[PASS]} or \texttt{[FAIL]}; in this report, a passed check is marked with~\pass{} and a failed one with~\fail.
\end{notebox}
